\documentclass[12pt,a4paper]{ctexart}
\usepackage{geometry}\geometry{margin=2.2cm}
\usepackage{graphicx,float,fancyhdr,hyperref,longtable,booktabs,enumitem,xcolor}
\pagestyle{fancy}\fancyhf{}\fancyhead[L]{IO 服务器 DLAS 本体}\fancyhead[R]{V4.0}\fancyfoot[C]{\thepage}
\title{\textbf{IO 服务器逆向工程本体\\DLAS 四层模型 v4.0}}
\author{约翰·刘}
\date{2026年7月13日}

\begin{document}\maketitle
\begin{abstract}
基于三周逆向工程数据（pcapng 95K报文、Oracle 61K测点、PDB 4732符号、199文件全量精读），
采用 DLAS (Data-Logic-Action-Security) 四层框架建立 IO 服务器知识本体。
覆盖 12 种设备类型、20 个真实设备 ID、8 个 IOMan 实例、5 条 OPC DA 访问路径。
\end{abstract}
\renewcommand{\abstractname}{摘要}

\section{Step 1: 盘点到齐——实体发现}

\subsection{Data 层: 服务器与网络}
\begin{table}[H]\centering
\begin{tabular}{lll}
\toprule
\textbf{实体} & \textbf{数量} & \textbf{标识} \\
\midrule
IO 服务器 & 1台 & 11.66.12.131 (Win2016, E:\textbackslash IO ServerOnLine) \\
pSpace 服务器 & 1台 & 11.66.12.130 (psAPI :8889) \\
Oracle 11g & 1台 & 11.66.12.129:1521/orcl (DQYTPROD) \\
RTU (Modbus) & 191台 & 11.248-250.x, 主动连接 :53001 \\
OPC DA 端点 & 5台 & 172.23.9.3/.23, 172.26.6.3, 172.23.18.194, 172.28.5.200 \\
OPC Server & Kepware 4.x & CLSID: 6E6170F0-FF2D-11D2-8087-00105AA8F840 \\
\bottomrule
\end{tabular}
\end{table}

\subsection{Data 层: 软件模块}
\begin{table}[H]\centering
\begin{tabular}{lrl}
\toprule
\textbf{模块} & \textbf{大小} & \textbf{职责} \\
\midrule
CommBridge.exe & 155KB & MFC主框架, TCP网关(191 RTU) \\
GPRSDLL.dll & 1.38MB & GPRS/CDMA协议栈 \\
psAPISDK.dll & 3.5MB & pSpace SDK (3525导出) \\
IOMan.exe & 299KB & IO管理器 (8实例) \\
IoMonitor.exe & 487KB & GUI监视 + Oracle写库 \\
IoProject.exe & 228KB & pSpace核心调度 \\
IoCommit.exe & 249KB & 批量写Oracle ×7 \\
\bottomrule
\end{tabular}
\end{table}

\subsection{Data 层: 12种保护装置}
\begin{longtable}{llll}
\toprule
\textbf{编码} & \textbf{名称} & \textbf{通道} & \textbf{典型测点} \\
\midrule
0x00 & DSL-31A 断路器 & 20 & Ia,Ib,Ic,Ua,Ub,Uc,P,cos$\phi$ \\
0x10 & DST-31A 变压器差动 & 15 & Ia差动,Ib,Ic,Ua,Ub,Uc \\
0x20 & DBPA-31A 备用电源 & 13 & 电压,开关状态 \\
0x30 & DSB-31A 变压器后备 & 20 & 同断路器 \\
0x40 & 电动机保护 & 19 & Ia,Ib,Ic,U,转速 \\
0x50 & DST-22D 变压器差动 & 20 & 同0x10 \\
0x60 & DSB-22D 变压器后备 & 20 & 同0x30 \\
0x70 & DSL-24D 断路器 & 20 & 同0x00 \\
0x80 & DGP-11 变压器差动保护 & 21 & 差动保护 \\
0x90 & DGP-12 变压器后备保护 & 24 & 后备保护 \\
0xA0 & DGP-13 接地保护 & 22 & 接地电流,零序 \\
0xB0 & DMP-31A 电动机保护 & 19 & 同0x40 \\
\bottomrule
\end{longtable}

\subsection{Data 层: 配置文件}
\begin{table}[H]\centering
\begin{tabular}{ll}
\toprule
\textbf{文件} & \textbf{内容} \\
\midrule
Device.ini & 12种保护装置 + ChangeData公式（标定8192） \\
IoChannelCfg.ini & CommBridge通道配置（DEV\_COUNT=3） \\
SqlFilSet.ini & Oracle连接 + SQL路径（EXECUTECYC=1000ms） \\
IoMonitor.ini & 提交时序（CommitRealSpan=300ms） \\
OPCClientCfg.ini & OPC Server ProgID→CLSID映射 \\
DeviceStruct.txt & OPC设备参数结构体(17字段) \\
DefinedStruct.txt & OPC点项结构体(15字段) \\
\bottomrule
\end{tabular}
\end{table}

\subsection{Data 层: 真实设备清单}
从 8 个 IOMan 实例命令行提取的 80 个设备 ID：

\begin{verbatim}
02012170058  02105100097  02105110008  02106290043  02106290052
02106290085  02107010048  02107030091  02107190091  02110080020
02110080028  02110110045  02110120089  02110150030  02110150041
02110150046  02110160086  02111260034  02111270046  02111270058
02204060100  02204060111  ... (共80台)
\end{verbatim}

\section{Step 2: 连线成网——关系矩阵}

\subsection{进程架构}
\begin{verbatim}
IoProject.exe (PID 5096)     ← pSpace Core, 总调度
    ├── CommBridge.exe (PID 19240)   ← TCP网关, 191 RTU
    │       └── DTU/ (16种DTU协议插件)
    ├── IOMan.exe ×5 (OPC_FC_Client) ← OPC DA DCOM → DCS
    ├── IOMan.exe ×1 (IM_A11_RTU)    ← Modbus RTU
    ├── IOMan.exe ×2 (A11)           ← A11 TCP → :8889
    ├── IoMonitor.exe (PID 18400)    ← GUI监视 + Oracle OLEDB
    └── IoCommit.exe ×7              ← 批量写Oracle
\end{verbatim}

\subsection{数据流链路}
\begin{verbatim}
RTU ──TCP:53001──→ CommBridge ──WM_COPYDATA──→ IoMonitor ──OLEDB──→ Oracle
DCS ──DCOM:135──→ OPC_FC_Client ──psAPI──→ pSpace ──IoCommit──→ Oracle
A11 ──TCP:8889───→ IM_A11_RTU ──psAPI─────→ pSpace ──IoCommit──→ Oracle
\end{verbatim}

\subsection{IPC 通信机制}
\begin{table}[H]\centering
\begin{tabular}{lll}
\toprule
\textbf{通道} & \textbf{技术} & \textbf{说明} \\
\midrule
IoProject→IoMonitor & FindWindow + WM\_COPYDATA & 窗口消息 \\
IoProject→IOMan & CreateProcess + 命令行 & 进程启动 \\
IOMan→IoMonitor & Global$\backslash$共享内存 + WM\_COPYDATA & 实时数据 \\
IoMonitor→IoCommit & Global$\backslash$共享内存 & 批量提交 \\
psAPISDK→pSpace & TCP :8889 + ACE框架 & 远程RPC \\
\bottomrule
\end{tabular}
\end{table}

\section{Step 3: 设卡立规——约束定义}

\subsection{时序约束}
\begin{table}[H]\centering
\begin{tabular}{lll}
\toprule
\textbf{约束} & \textbf{阈值} & \textbf{来源} \\
\midrule
RTU 轮询间隔 & 1000ms & IoMonitor \\
OPC 刷新时间 & 设备参数[164] & DeviceStruct \\
数据提交(实时) & 300ms & IoMonitor.ini \\
数据提交(历史) & 500ms & IoMonitor.ini \\
提交批次大小 & 15000点 & IoMonitor.ini \\
SQL 执行周期 & 1000ms & SqlFilSet.ini \\
CommBridge 超时 & 30s & IoChannelCfg.ini \\
dgiot TCP 超时 & 120s & commbridge\_server.py \\
\bottomrule
\end{tabular}
\end{table}

\subsection{数据转换约束 (Device.ini ChangeData)}
\begin{table}[H]\centering
\begin{tabular}{lll}
\toprule
\textbf{系数} & \textbf{值} & \textbf{公式} \\
\midrule
C0 & 170/8192 & Ia,Ib,Ic,Ua,Ub,Uc（电流/电压） \\
C1 & 8.5/8192 & Iac（接地电流） \\
C2 & 170/8192 & Uphase（相电压） \\
C3 & 170×8.5/8192 & P（有功功率） \\
C4 & 1/8192 & cos$\phi$（功率因数） \\
C5 & 2/8192 & F=50+Y×2/8192（频率） \\
C6-9 & 1 & 状态量（直通） \\
\bottomrule
\end{tabular}
\end{table}

\subsection{数据校验约束 (L1-L5)}
\begin{table}[H]\centering
\begin{tabular}{lll}
\toprule
\textbf{层级} & \textbf{检查项} & \textbf{判定标准} \\
\midrule
L1 & 帧Len匹配 & abs(expected$-$actual) $\leq$ 2 \\
L2 & 值范围 & 电流0-500A / 电压100-400V / 功率0-300kW \\
L3 & 三相平衡 & (I$_{max}$$-$I$_{min}$)/I$_{avg}$ $<$ 25\% \\
L4 & 历时一致性 & 相邻值变化 $<$ 50\% \\
L5 & Oracle对标 & |A$-$B|/A $<$ 1\% \\
\bottomrule
\end{tabular}
\end{table}

\subsection{Security 层: 安全约束}
\begin{table}[H]\centering
\begin{tabular}{lll}
\toprule
\textbf{系统} & \textbf{认证方式} & \textbf{凭据摘要} \\
\midrule
WinRM :5985 & NTLM & administrator / **** \\
Oracle :1521 & OLEDB & DQYTPROD / dqyta11\_**** \\
pSpace :8889 & psAPI & admin / DQYTA11\_**** \\
DCOM :135 & CoInitSecurity & DCS侧授权 \\
\bottomrule
\end{tabular}
\end{table}

\section{Data 层补充: pSpace 双命名空间}

逆向发现 pSpace 存在两套独立的 Tag 命名体系:

\begin{table}[H]\centering
\begin{tabular}{lll}
\toprule
\textbf{属性} & \textbf{层级路径} & \textbf{数字 ID} \\
\midrule
设备类型 & 油田泵油单元 & 保护装置 \\
Tag 命名 & /CY1C7K/G138VS405/JD... & 5000, 5500, 7500, ... \\
访问方式 & psAPI Tag Query/Path & psAPI Real ReadList(ID) \\
缓存文件 & TagID\_CY*.dat (30 文件) & 动态解析 \\
设备 ID 格式 & G138VS405, JD010707... & 02204060100, 02105100097 \\
数据量 & 数千 (pPluse.dat 123KB) & 14 块 \\
发现来源 & IOMan PDB + TagID 文件 & pSpace ReadList 扫描 \\
\bottomrule
\end{tabular}
\end{table}

\section{Step 4: 闭环验证}

\begin{table}[H]\centering
\begin{tabular}{lll}
\toprule
\textbf{验证项} & \textbf{方法} & \textbf{状态} \\
\midrule
CommBridge协议帧 & 95K报文全量 & 异常率0.017\% \\
RTU注册包格式 & 2211包, 413种ID & 全部0xAA+ASCII+0x0D \\
float32/int16识别 & ByteCount\%4判定 & 67\%/33\% \\
数据公式验证 & 170/8192标定 & CT变比30:1→PF=0.77 \\
Oracle点位映射 & SYS\_POINTRELATION→TAGPAR & 54口井×23通道 \\
OPC设备配置 & DeviceStruct反编译 & 17字段验证 \\
SQL执行链路 & SqlFilSet.ini & EXECUTECYC 1000ms→Oracle \\
崩溃历史 & 13个dump文件 & AdoInterface×8, CommBridge×5 \\
dgiot Server替代 & 真实注册包回放 & 查询帧格式一致 \\
psAPISDK直连 & ctypes + ReadList & Tag 5000=1.875 GOOD \\
\bottomrule
\end{tabular}
\end{table}

\section{OPC DA 五条访问路径}

\begin{longtable}{lllcc}
\toprule
\textbf{\#} & \textbf{路径} & \textbf{访问层} & \textbf{实时} & \textbf{状态} \\
\midrule
① & Python Mock 全栈模拟 & 本地socket & -- & 已验证(56设备) \\
② & dgiot\_lite Oracle管线 & SQL同步 & 否 & 已通(26万条) \\
③ & psAPISDK直连pSpace & TCP :8889 & 是 & \textbf{已通(自主读Tag)} \\
④ & IoProject注入自定义IOMan & 进程注入 & 是 & 格式已验证 \\
⑤ & DCOM直连Kepware OPC DA & WinRM中转 & 是 & 脚本就绪 \\
\bottomrule
\end{longtable}

\section{实体关系总图}

\begin{verbatim}
┌─────────────────────────────────────────────────────┐
│                IO 服务器 DLAS 本体                   │
│                                                     │
│  Data 层                                            │
│    12种设备·20设备ID·7服务·3协议·10系数·199文件        │
│         │                                           │
│         ├──[监测]──→ Logic 层                        │
│         │   时序约束·数据校验L1-L5·12通道映射           │
│         │       │                                   │
│         │       ├──[触发]──→ Action 层               │
│         │       │   RTU→CB→IoMon→IoCom→Oracle       │
│         │       │   五条OPC DA访问路径                │
│         │       │       │                           │
│         │       │       └──[审计]──→ Security 层      │
│         │       │           凭据·权限·加密·合规         │
│         │       │                                   │
│         └──[闭环反馈]──→ 新一轮实体发现                │
└─────────────────────────────────────────────────────┘
\end{verbatim}

\section{版本历史}
\begin{table}[H]\centering
\begin{tabular}{lll}
\toprule
\textbf{版本} & \textbf{日期} & \textbf{更新} \\
\midrule
v1.0 & 2026-07-10 & 初始版本, 基于pcapng+Oracle+PDB \\
v2.0 & 2026-07-11 & 新增DeviceStruct反编译+SqlFilSet+崩溃dump \\
v3.0 & 2026-07-12 & 新增CommBridge真实协议+数据校验L1-L5 \\
v4.0 & 2026-07-13 & IOMan命令行格式+IPC架构+psAPISDK直连+五条路径 \\
\bottomrule
\end{tabular}
\end{table}

\end{document}
